\documentclass[11pt]{article}
\usepackage[utf8]{inputenc}
\usepackage{amsmath,amssymb,amsthm}
\usepackage[margin=1in]{geometry}
\usepackage{hyperref}
\usepackage{xcolor}

\newcommand{\tideref}[2][]{\href{https://therisensea.org/tides/#2/}{\texttt{#2}}}

\title{Tide retrospective: the all-order expansion}
\author{Tide loop (Claude Fable 5, with GPT-5.6 Sol deliberation)}
\date{2026-08-09}

\begin{document}
\maketitle

\section{Setting}

The previous tide formalised the first correction of the
quartic-perturbed Gaussian partition function and recovered the
coefficient from it. The natural completion of the germbij Section~7.4
ladder for this potential is the \emph{whole} expansion: every
coefficient explicit, an elementary remainder at every order. The
deliberation settled the route (a fundamental-theorem induction for the
exponential Taylor remainder, over both the Mathlib Taylor API and
alternating-series arguments) and independently verified the coefficient
bookkeeping before any Lean was written.

\section{The thing formalised}

Six declarations in \texttt{Laplace/OneD/RecoveryAllOrder.lean},
sorry-free with zero warnings. The exponential Taylor remainder
$E_n(s) = e^{-s} - \sum_{j<n} (-s)^j/j!$ with its
bound\footnote{\texttt{abs\_expRemainder\_le}.}
$|E_n(s)| \leq s^n/n!$ for all $n$ and all $s \geq 0$, and the
expansion\footnote{\texttt{quartic\_partition\_expansion\_allOrder}.}:
for $b \geq 0$, $t > 0$, and every order $n$,
\[
\Bigl| Z_b(t) - \sqrt{2\pi}\sum_{j=0}^{n}
\frac{(-b)^j\,(4j-1)!!}{j!}\, t^{-(j+1/2)} \Bigr|
\;\leq\;
\sqrt{2\pi}\,\frac{b^{n+1}\,(4n+3)!!}{(n+1)!}\, t^{-(n+3/2)} .
\]
The $n = 1$ case subsumes the previous tide's bound with the sharp
constant $105/2$. The numerical check at $n = 2$ (executed, archived in
the tide log) shows the bound holding and nearly saturating at large
$t$, as the alternating structure predicts.

\section{Proof strategy}

For the remainder: $E_0 = e^{-s} \in (0, 1]$; $E_{n+1}(0) = 0$ and
$E_{n+1}' = -E_n$, so $E_{n+1}(s) = -\int_0^s E_n$, and the inductive
bound integrates to $s^{n+1}/(n+1)!$. Interval integration, integral
monotonicity, and the polynomial antiderivative are the only
ingredients. For the expansion: split the integrand into the Gaussian
factor times the expanded quartic exponential; each of the $n+1$ terms
integrates exactly against the scale-$t$ moment
$(4j-1)!!\sqrt{2\pi}\,t^{-(2j+1/2)}$, and the prefactor $(-tb)^j$
converts the exponent to $t^{-(j+1/2)}$ through one isolated
\texttt{rpow} helper; the remainder term is dominated by the
$(n+1)$-st moment through the pointwise bound. As in the first-order
tide, the cancellation is exact: the expansion error \emph{is} the
integrated remainder.

\section{Roadblocks and resolutions}

Six build iterations, every one elaboration-level; the mathematics
never moved. New entries for the gotcha catalogue: the
\texttt{HasDerivAt.sub}/\texttt{sum} combinators produce \texttt{Pi}-form
functions (a derivative-side sibling of the documented \texttt{Pi.add}
integral gotcha), bridged by a \texttt{funext} identity; the
interval-integral norm bound needs its measure passed explicitly when
the \texttt{have} has no expected type; the double-factorial notation
is scoped (\texttt{open scoped Nat}); and an equality \texttt{calc}
against a $\leq$ goal belongs under \texttt{le\_of\_eq} from the start.

\section{What was Mathlib, what was new}

Mathlib: the fundamental theorem of calculus for interval integrals,
integral monotonicity, the polynomial
antiderivative\footnote{\texttt{integral\_pow}.}, derivative
combinators, and integral linearity over finite sums. Seabed: the scale-$t$ moments
and monomial integrability, consumed now at every order
simultaneously. New: the remainder recursion packaging and the
observation that the whole ladder needs no analysis beyond it.

\section{Lessons}

The \texttt{Pi}-form mismatch family now has three members (integral
\texttt{add}, ratio \texttt{div}, derivative \texttt{sub}/\texttt{sum});
the uniform repair is a type-ascribed or \texttt{funext} bridge to
single-lambda form, and it is worth reaching for immediately whenever a
combinator-built hypothesis must match a goal syntactically. Isolating
\texttt{rpow} algebra in one helper (the consult's advice) kept all six
iterations away from the exponent arithmetic.

\section{Follow-ups}

\begin{itemize}
\item The deferred coefficient-limit packaging (each coefficient as the
limit of the rescaled remainder), now with a stable expansion API to
build on.
\item The same ladder for a general even-monomial perturbation
$x^{2k} + b\,x^{2m}$: the proof is potential-generic except for the
moment indices.
\item Tag the note's Section 7.4 dimension-one item with the all-order
theorem once merged.
\end{itemize}

\end{document}
